\section{Bố cục luận văn}
\label{sec:organization}

Luận văn được tổ chức thành năm chương. \textbf{Chương 1} giới thiệu bối cảnh nghiên cứu về thực khuẩn thể, các thách thức trong phân loại lối sống phage từ dữ liệu metagenomic, mục tiêu và đóng góp chính, phạm vi nghiên cứu, và bố cục luận văn. \textbf{Chương 2} cung cấp nền tảng lý thuyết về sinh học thực khuẩn thể, metagenomics, các kiến trúc học sâu (CNN, Transformer, BERT), mô hình DNABERT-2, và tổng quan các công trình liên quan về phân loại lối sống phage. \textbf{Chương 3} trình bày chi tiết phương pháp PhaBERT-CNN, bao gồm quy trình xây dựng dữ liệu (sliding window, reverse complement augmentation, Random Undersampling), kiến trúc mô hình (DNABERT-2 backbone, Modified TextCNN với nhánh CNN đa tỷ lệ và nhánh attention pooling, đầu phân loại), và chiến lược huấn luyện hai giai đoạn với discriminative learning rates. \textbf{Chương 4} trình bày thiết lập thực nghiệm, kết quả so sánh PhaBERT-CNN với DNABERT-2 baseline và các phương pháp tiên tiến (PhaTYP, DeePhage, ProkBERT) trên bốn nhóm độ dài contig. \textbf{Chương 5} tổng kết đóng góp của luận văn, phân tích ưu điểm và hạn chế của phương pháp dựa trên kết quả thực nghiệm, và đề xuất các hướng nghiên cứu tương lai.
Luận văn được tổ chức thành năm chương với cấu trúc logic như sau:

\textbf{Chương 1: Giới thiệu}

Chương này giới thiệu bối cảnh nghiên cứu về thực khuẩn thể và tầm quan trọng của việc phân loại lối sống phage trong nghiên cứu vi sinh vật học và ứng dụng y sinh học. Chương trình bày các thách thức chính trong phân loại lối sống phage từ dữ liệu metagenomic, bao gồm hạn chế về dữ liệu, thiếu gen đánh dấu phổ quát, độ dài đoạn DNA ngắn, và vấn đề tương đồng trình tự. Chương cũng phát biểu rõ ràng bài toán nghiên cứu, mục tiêu, đóng góp chính, và phạm vi của luận văn.

\textbf{Chương 2: Cơ sở lý thuyết}

Chương này cung cấp nền tảng lý thuyết toàn diện cần thiết để hiểu phương pháp đề xuất. Chương bắt đầu với kiến thức sinh học về thực khuẩn thể, bao gồm phân loại, chu trình sống, và vai trò sinh thái. Tiếp theo, chương giới thiệu các khái niệm về metagenomics và giải trình tự thế hệ mới, giải thích cách các contig được tạo ra từ dữ liệu metagenomic. Phần tiếp theo trình bày các nền tảng học máy và học sâu, bao gồm mạng nơ-ron tích chập (CNN), kiến trúc transformer, và cơ chế attention. Chương dành phần quan trọng để giải thích chi tiết về BERT và các biến thể của nó cho dữ liệu genome (DNABERT, DNABERT-2), bao gồm masked language modeling, tokenization strategies, và transfer learning. Cuối cùng, chương tổng hợp các công trình liên quan về phân loại lối sống phage, phân tích ưu nhược điểm của các phương pháp học máy truyền thống (PHACTS, PhageAI, BACPHLIP) và các phương pháp học sâu hiện đại (DeePhage, PhaTYP, DeepPL, ProkBERT).

\textbf{Chương 3: Phương pháp đề xuất}

Chương này trình bày chi tiết phương pháp PhaBERT-CNN được đề xuất. Chương bắt đầu với tổng quan về kiến trúc tổng thể và động lực thiết kế. Tiếp theo, chương mô tả quy trình xây dựng tập dữ liệu, bao gồm nguồn dữ liệu, phương pháp tạo contig bằng sliding window, data augmentation thông qua reverse complement, và chiến lược xử lý mất cân bằng lớp. Chương sau đó giải thích chi tiết về mô hình nền tảng DNABERT-2, bao gồm chiến lược tokenization BPE, các cải tiến kiến trúc (ALiBi, Flash Attention), và corpus tiền huấn luyện đa loài. Phần cốt lõi của chương mô tả kiến trúc PhaBERT-CNN với hai nhánh trích xuất đặc trưng song song: (1) multi-scale CNN feature extractor với ba kernel sizes khác nhau, và (2) attention-based global feature extractor. Chương kết thúc bằng việc trình bày chiến lược huấn luyện hai giai đoạn với learning rate phân biệt, bao gồm các chi tiết về optimizer, learning rate schedule, và early stopping.

\textbf{Chương 4: Thực nghiệm và đánh giá}

Chương này trình bày thiết lập thí nghiệm toàn diện và kết quả đánh giá. Chương bắt đầu với mô tả chi tiết về môi trường thực nghiệm, bao gồm cấu hình phần cứng, phần mềm, và các hyperparameters. Tiếp theo, chương giải thích các metrics đánh giá (sensitivity, specificity, accuracy) và phương pháp validation (stratified 5-fold cross-validation). Chương trình bày kết quả thực nghiệm theo hai phần chính: (1) so sánh PhaBERT-CNN với DNABERT-2 baseline để đánh giá tác động của các thành phần CNN và attention, và (2) so sánh với các phương pháp tiên tiến (PhaTYP, DeePhage, ProkBERT) trên bốn nhóm độ dài contig. Kết quả được trình bày thông qua các bảng số liệu chi tiết và biểu đồ trực quan. Chương cũng bao gồm phân tích chi tiết về hiệu suất trên từng nhóm độ dài, phân tích lỗi, và thảo luận về các trường hợp mô hình hoạt động tốt và các trường hợp thất bại. Cuối cùng, chương thảo luận về trade-off giữa độ chính xác và thời gian tính toán, cũng như các hạn chế của phương pháp đề xuất.

\textbf{Chương 5: Kết luận}

Chương cuối cùng tổng kết các kết quả chính của luận văn và đóng góp của nghiên cứu. Chương nhắc lại bài toán nghiên cứu và tóm tắt cách PhaBERT-CNN giải quyết các thách thức đã xác định. Chương tổng hợp các phát hiện quan trọng từ thực nghiệm, bao gồm hiệu suất vượt trội trên các contig ngắn và trung bình, sự cân bằng giữa sensitivity và specificity, và so sánh với các phương pháp hiện có. Chương thảo luận về các hạn chế của nghiên cứu, bao gồm phụ thuộc vào chất lượng dữ liệu đầu vào, nhu cầu về tài nguyên tính toán, và các giả định về dữ liệu. Cuối cùng, chương đề xuất các hướng nghiên cứu tương lai, bao gồm: (1) đánh giá trên dữ liệu metagenomic thực tế từ các môi trường đa dạng, (2) khảo sát các mô hình nền tảng genome khác, (3) phát triển các chiến lược xử lý mất cân bằng lớp hiệu quả hơn, (4) tối ưu hóa cho triển khai thực tế, và (5) mở rộng sang các nhiệm vụ phân loại phage khác như phân loại họ phage hoặc dự đoán vi khuẩn chủ.

\vspace{1em}

Cấu trúc này được thiết kế để dẫn dắt người đọc từ động lực và nền tảng lý thuyết, qua phương pháp đề xuất chi tiết, đến đánh giá thực nghiệm toàn diện, và cuối cùng là kết luận và hướng phát triển. Mỗi chương xây dựng dựa trên các chương trước, tạo thành một câu chuyện nghiên cứu mạch lạc và logic.
